
\documentclass{article}
\usepackage{colortbl}
\usepackage{makecell}
\usepackage{multirow}
\usepackage{supertabular}

\begin{document}

\newcounter{utterance}

\centering \large Interaction Transcript for game `dond', experiment `semi\_it', episode 14 with claude{-}sonnet{-}4{-}20250514{-}no{-}reasoning{-}t1.0.
\vspace{24pt}

{ \footnotesize  \setcounter{utterance}{1}
\setlength{\tabcolsep}{0pt}
\begin{supertabular}{c@{$\;$}|p{.15\linewidth}@{}p{.15\linewidth}p{.15\linewidth}p{.15\linewidth}p{.15\linewidth}p{.15\linewidth}}
    \# & $\;$A & \multicolumn{4}{c}{Game Master} & $\;\:$B\\
    \hline

    \theutterance \stepcounter{utterance}  
    & & \multicolumn{4}{p{0.6\linewidth}}{
        \cellcolor[rgb]{0.9,0.9,0.9}{
            \makecell[{{p{\linewidth}}}]{
                \texttt{\tiny{[P1$\langle$GM]}}
                \texttt{State giocando a un gioco di negoziazione in cui dovete accordarvi su come dividere una serie di oggetti tra voi e un altro giocatore.} \\
\\ 
\texttt{Regole:} \\
\texttt{(a) A Lei e all'altro giocatore viene dato un insieme di oggetti. Ognuno di voi riceve anche una funzione di valore segreta, che rappresenta il valore di ciascun tipo di oggetto.} \\
\texttt{(b) Si scambiano messaggi con l'altro giocatore per concordare chi si aggiudica gli oggetti. Potete inviare un massimo di 5 messaggi ciascuno, oppure terminare in anticipo facendo una proposta segreta in qualsiasi momento.} \\
\texttt{(c) A ciascuno di voi viene chiesto di inviare una proposta segreta indicando gli oggetti che desiderate, formattata tra parentesi quadre come segue: "{[}Proposta: <numero> <nome oggetto>, <numero> <nome oggetto>, <...>{]}".} \\
\texttt{(d) Se le vostre proposte sono complementari, cioè ci sono abbastanza oggetti per soddisfare entrambe le proposte, a ciascun giocatore viene assegnato un punteggio basato sulla somma dei valori degli oggetti ricevuti. In caso contrario, entrambi ricevono zero punti.} \\
\texttt{(e) L'obiettivo è massimizzare il proprio punteggio.} \\
\\ 
\texttt{Cominciamo.} \\
\\ 
\texttt{L'insieme degli oggetti disponibili è:} \\
\\ 
\texttt{1 cornice, 1 scarpa, 3 martelli} \\
\\ 
\texttt{La funzione valore segreta è:} \\
\\ 
\texttt{cornice: 8, scarpa: 2, martello: 0} \\
\\ 
\texttt{IMPORTANTE: i vostri messaggi, a meno che non si tratti di una proposta segreta, vengono trasmessi direttamente all'altro giocatore, quindi non includete alcuna risposta alle regole o testo di annuncio del vostro messaggio. Per fare una proposta segreta, utilizzate il formato indicato. Non utilizzare delle parentesi quadre quando si comunica all'altro giocatore, altrimenti verrà interpretata come una proposta segreta.} \\
            }
        }
    }
    & \\ \\

    \theutterance \stepcounter{utterance}  
    & \multicolumn{4}{p{0.6\linewidth}}{
        \cellcolor[rgb]{0.9,0.9,0.9}{
            \makecell[{{p{\linewidth}}}]{
                \texttt{\tiny{[P1$\rangle$GM]}}
                \texttt{Ciao! Vedo che abbiamo 1 cornice, 1 scarpa e 3 martelli da dividere. Sono interessato principalmente alla cornice. Che tipo di oggetti ti interessano di più? Forse possiamo trovare un accordo che funzioni per entrambi.} \\
            }
        }
    }
    & & \\ \\

    \theutterance \stepcounter{utterance}  
    & & & \multicolumn{2}{p{0.3\linewidth}}{
        \cellcolor[rgb]{0.9,0.9,0.9}{
            \makecell[{{p{\linewidth}}}]{
                \texttt{\tiny{[GM$|$GM]}}
                \texttt{continue} \\
            }
        }
    }
    & & \\ \\

    \theutterance \stepcounter{utterance}  
    & & & \multicolumn{2}{p{0.3\linewidth}}{
        \cellcolor[rgb]{0.9,0.9,0.9}{
            \makecell[{{p{\linewidth}}}]{
                \texttt{\tiny{[GM$|$GM]}}
                \texttt{Ciao! Vedo che abbiamo 1 cornice, 1 scarpa e 3 martelli da dividere. Sono interessato principalmente alla cornice. Che tipo di oggetti ti interessano di più? Forse possiamo trovare un accordo che funzioni per entrambi.} \\
            }
        }
    }
    & & \\ \\

    \theutterance \stepcounter{utterance}  
    & & \multicolumn{4}{p{0.6\linewidth}}{
        \cellcolor[rgb]{0.9,0.9,0.9}{
            \makecell[{{p{\linewidth}}}]{
                \texttt{\tiny{[GM$\rangle$P2]}}
                \texttt{State giocando a un gioco di negoziazione in cui dovete accordarvi su come dividere una serie di oggetti tra voi e un altro giocatore.} \\
\\ 
\texttt{Regole:} \\
\texttt{(a) A Lei e all'altro giocatore viene dato un insieme di oggetti. Ognuno di voi riceve anche una funzione di valore segreta, che rappresenta il valore di ciascun tipo di oggetto.} \\
\texttt{(b) Si scambiano messaggi con l'altro giocatore per concordare chi si aggiudica gli oggetti. Potete inviare un massimo di 5 messaggi ciascuno, oppure terminare in anticipo facendo una proposta segreta in qualsiasi momento.} \\
\texttt{(c) A ciascuno di voi viene chiesto di inviare una proposta segreta indicando gli oggetti che desiderate, formattata tra parentesi quadre come segue: "{[}Proposta: <numero> <nome oggetto>, <numero> <nome oggetto>, <...>{]}".} \\
\texttt{(d) Se le vostre proposte sono complementari, cioè ci sono abbastanza oggetti per soddisfare entrambe le proposte, a ciascun giocatore viene assegnato un punteggio basato sulla somma dei valori degli oggetti ricevuti. In caso contrario, entrambi ricevono zero punti.} \\
\texttt{(e) L'obiettivo è massimizzare il proprio punteggio.} \\
\\ 
\texttt{Cominciamo.} \\
\\ 
\texttt{L'insieme degli oggetti disponibili è:} \\
\\ 
\texttt{1 cornice, 1 scarpa, 3 martelli} \\
\\ 
\texttt{La funzione valore segreta è:} \\
\\ 
\texttt{cornice: 6, scarpa: 1, martello: 1} \\
\\ 
\texttt{IMPORTANTE: i vostri messaggi, a meno che non si tratti di una proposta segreta, vengono trasmessi direttamente all'altro giocatore, quindi non includete alcuna risposta alle regole o testo di annuncio del vostro messaggio. Per fare una proposta segreta, utilizzate il formato indicato. Non utilizzare delle parentesi quadre quando si comunica all'altro giocatore, altrimenti verrà interpretata come una proposta segreta.} \\
\\ 
\\ 
\texttt{Ciao! Vedo che abbiamo 1 cornice, 1 scarpa e 3 martelli da dividere. Sono interessato principalmente alla cornice. Che tipo di oggetti ti interessano di più? Forse possiamo trovare un accordo che funzioni per entrambi.} \\
            }
        }
    }
    & \\ \\

    \theutterance \stepcounter{utterance}  
    & & & \multicolumn{4}{p{0.6\linewidth}}{
        \cellcolor[rgb]{0.9,0.9,0.9}{
            \makecell[{{p{\linewidth}}}]{
                \texttt{\tiny{[GM$\langle$P2]}}
                \texttt{Ciao! Grazie per aver iniziato la conversazione. Anch'io sono molto interessato alla cornice {-} è l'oggetto che mi attira di più. Per quanto riguarda gli altri oggetti, non ho preferenze particolari tra scarpe e martelli. } \\
\\ 
\texttt{Visto che siamo entrambi interessati alla cornice, come pensi che potremmo risolvere questa situazione? Hai qualche idea per un accordo equo?} \\
            }
        }
    }
     \\ \\

    \theutterance \stepcounter{utterance}  
    & & & \multicolumn{2}{p{0.3\linewidth}}{
        \cellcolor[rgb]{0.9,0.9,0.9}{
            \makecell[{{p{\linewidth}}}]{
                \texttt{\tiny{[GM$|$GM]}}
                \texttt{continue} \\
            }
        }
    }
    & & \\ \\

    \theutterance \stepcounter{utterance}  
    & & & \multicolumn{2}{p{0.3\linewidth}}{
        \cellcolor[rgb]{0.9,0.9,0.9}{
            \makecell[{{p{\linewidth}}}]{
                \texttt{\tiny{[GM$|$GM]}}
                \texttt{Ciao! Grazie per aver iniziato la conversazione. Anch'io sono molto interessato alla cornice {-} è l'oggetto che mi attira di più. Per quanto riguarda gli altri oggetti, non ho preferenze particolari tra scarpe e martelli. } \\
\\ 
\texttt{Visto che siamo entrambi interessati alla cornice, come pensi che potremmo risolvere questa situazione? Hai qualche idea per un accordo equo?} \\
            }
        }
    }
    & & \\ \\

    \theutterance \stepcounter{utterance}  
    & & \multicolumn{4}{p{0.6\linewidth}}{
        \cellcolor[rgb]{0.9,0.9,0.9}{
            \makecell[{{p{\linewidth}}}]{
                \texttt{\tiny{[P1$\langle$GM]}}
                \texttt{Ciao! Grazie per aver iniziato la conversazione. Anch'io sono molto interessato alla cornice {-} è l'oggetto che mi attira di più. Per quanto riguarda gli altri oggetti, non ho preferenze particolari tra scarpe e martelli. } \\
\\ 
\texttt{Visto che siamo entrambi interessati alla cornice, come pensi che potremmo risolvere questa situazione? Hai qualche idea per un accordo equo?} \\
            }
        }
    }
    & \\ \\

    \theutterance \stepcounter{utterance}  
    & \multicolumn{4}{p{0.6\linewidth}}{
        \cellcolor[rgb]{0.9,0.9,0.9}{
            \makecell[{{p{\linewidth}}}]{
                \texttt{\tiny{[P1$\rangle$GM]}}
                \texttt{Capisco che siamo entrambi interessati alla cornice {-} questa è sicuramente la questione principale da risolvere! } \\
\\ 
\texttt{Dato che c'è solo una cornice e la vogliamo entrambi, potrei essere disposto a rinunciarvi in cambio di tutti gli altri oggetti disponibili (la scarpa e tutti e 3 i martelli). In questo modo tu otterresti l'oggetto che desideri di più, mentre io avrei una compensazione ragionevole con tutto il resto.} \\
\\ 
\texttt{Cosa ne pensi di questa divisione: tu prendi la cornice, io prendo la scarpa e i 3 martelli?} \\
            }
        }
    }
    & & \\ \\

    \theutterance \stepcounter{utterance}  
    & & & \multicolumn{2}{p{0.3\linewidth}}{
        \cellcolor[rgb]{0.9,0.9,0.9}{
            \makecell[{{p{\linewidth}}}]{
                \texttt{\tiny{[GM$|$GM]}}
                \texttt{continue} \\
            }
        }
    }
    & & \\ \\

    \theutterance \stepcounter{utterance}  
    & & & \multicolumn{2}{p{0.3\linewidth}}{
        \cellcolor[rgb]{0.9,0.9,0.9}{
            \makecell[{{p{\linewidth}}}]{
                \texttt{\tiny{[GM$|$GM]}}
                \texttt{Capisco che siamo entrambi interessati alla cornice {-} questa è sicuramente la questione principale da risolvere! } \\
\\ 
\texttt{Dato che c'è solo una cornice e la vogliamo entrambi, potrei essere disposto a rinunciarvi in cambio di tutti gli altri oggetti disponibili (la scarpa e tutti e 3 i martelli). In questo modo tu otterresti l'oggetto che desideri di più, mentre io avrei una compensazione ragionevole con tutto il resto.} \\
\\ 
\texttt{Cosa ne pensi di questa divisione: tu prendi la cornice, io prendo la scarpa e i 3 martelli?} \\
            }
        }
    }
    & & \\ \\

    \theutterance \stepcounter{utterance}  
    & & \multicolumn{4}{p{0.6\linewidth}}{
        \cellcolor[rgb]{0.9,0.9,0.9}{
            \makecell[{{p{\linewidth}}}]{
                \texttt{\tiny{[GM$\rangle$P2]}}
                \texttt{Capisco che siamo entrambi interessati alla cornice {-} questa è sicuramente la questione principale da risolvere! } \\
\\ 
\texttt{Dato che c'è solo una cornice e la vogliamo entrambi, potrei essere disposto a rinunciarvi in cambio di tutti gli altri oggetti disponibili (la scarpa e tutti e 3 i martelli). In questo modo tu otterresti l'oggetto che desideri di più, mentre io avrei una compensazione ragionevole con tutto il resto.} \\
\\ 
\texttt{Cosa ne pensi di questa divisione: tu prendi la cornice, io prendo la scarpa e i 3 martelli?} \\
            }
        }
    }
    & \\ \\

    \theutterance \stepcounter{utterance}  
    & & & \multicolumn{4}{p{0.6\linewidth}}{
        \cellcolor[rgb]{0.9,0.9,0.9}{
            \makecell[{{p{\linewidth}}}]{
                \texttt{\tiny{[GM$\langle$P2]}}
                \texttt{Hmm, apprezzo la tua proposta e capisco il ragionamento. Tuttavia, non sono sicuro che quella divisione sia del tutto equa per me. La cornice è sicuramente l'oggetto principale, ma rinunciare a tutto il resto (4 oggetti in totale) per avere solo la cornice mi sembra un po' sbilanciato.} \\
\\ 
\texttt{Che ne dici se invece io prendo la cornice e 1 martello, mentre tu prendi la scarpa e 2 martelli? In questo modo entrambi otteniamo l'oggetto che ci interessa di più (per me la cornice, per te una buona quantità degli altri oggetti) ma la divisione è un po' più equilibrata.} \\
\\ 
\texttt{Come ti suona questa alternativa?} \\
            }
        }
    }
     \\ \\

    \theutterance \stepcounter{utterance}  
    & & & \multicolumn{2}{p{0.3\linewidth}}{
        \cellcolor[rgb]{0.9,0.9,0.9}{
            \makecell[{{p{\linewidth}}}]{
                \texttt{\tiny{[GM$|$GM]}}
                \texttt{continue} \\
            }
        }
    }
    & & \\ \\

    \theutterance \stepcounter{utterance}  
    & & & \multicolumn{2}{p{0.3\linewidth}}{
        \cellcolor[rgb]{0.9,0.9,0.9}{
            \makecell[{{p{\linewidth}}}]{
                \texttt{\tiny{[GM$|$GM]}}
                \texttt{Hmm, apprezzo la tua proposta e capisco il ragionamento. Tuttavia, non sono sicuro che quella divisione sia del tutto equa per me. La cornice è sicuramente l'oggetto principale, ma rinunciare a tutto il resto (4 oggetti in totale) per avere solo la cornice mi sembra un po' sbilanciato.} \\
\\ 
\texttt{Che ne dici se invece io prendo la cornice e 1 martello, mentre tu prendi la scarpa e 2 martelli? In questo modo entrambi otteniamo l'oggetto che ci interessa di più (per me la cornice, per te una buona quantità degli altri oggetti) ma la divisione è un po' più equilibrata.} \\
\\ 
\texttt{Come ti suona questa alternativa?} \\
            }
        }
    }
    & & \\ \\

    \theutterance \stepcounter{utterance}  
    & & \multicolumn{4}{p{0.6\linewidth}}{
        \cellcolor[rgb]{0.9,0.9,0.9}{
            \makecell[{{p{\linewidth}}}]{
                \texttt{\tiny{[P1$\langle$GM]}}
                \texttt{Hmm, apprezzo la tua proposta e capisco il ragionamento. Tuttavia, non sono sicuro che quella divisione sia del tutto equa per me. La cornice è sicuramente l'oggetto principale, ma rinunciare a tutto il resto (4 oggetti in totale) per avere solo la cornice mi sembra un po' sbilanciato.} \\
\\ 
\texttt{Che ne dici se invece io prendo la cornice e 1 martello, mentre tu prendi la scarpa e 2 martelli? In questo modo entrambi otteniamo l'oggetto che ci interessa di più (per me la cornice, per te una buona quantità degli altri oggetti) ma la divisione è un po' più equilibrata.} \\
\\ 
\texttt{Come ti suona questa alternativa?} \\
            }
        }
    }
    & \\ \\

    \theutterance \stepcounter{utterance}  
    & \multicolumn{4}{p{0.6\linewidth}}{
        \cellcolor[rgb]{0.9,0.9,0.9}{
            \makecell[{{p{\linewidth}}}]{
                \texttt{\tiny{[P1$\rangle$GM]}}
                \texttt{Capisco il tuo punto di vista sulla necessità di un accordo più equilibrato. La tua proposta è ragionevole {-} tu avresti la cornice e 1 martello, mentre io avrei la scarpa e 2 martelli.} \\
\\ 
\texttt{Tuttavia, dal mio punto di vista, rinunciare alla cornice (che è chiaramente l'oggetto più desiderabile per entrambi) merita una compensazione un po' più sostanziosa. } \\
\\ 
\texttt{Che ne dici di un compromesso intermedio? Tu prendi la cornice, io prendo la scarpa e 2 martelli, e tu tieni 1 martello {-} che è esattamente quello che hai proposto tu! } \\
\\ 
\texttt{Scusa, rileggendo meglio, hai ragione che la tua proposta era già un compromesso equo. Sono d'accordo con la tua divisione: tu cornice + 1 martello, io scarpa + 2 martelli. Mi sembra un accordo giusto per entrambi.} \\
            }
        }
    }
    & & \\ \\

    \theutterance \stepcounter{utterance}  
    & & & \multicolumn{2}{p{0.3\linewidth}}{
        \cellcolor[rgb]{0.9,0.9,0.9}{
            \makecell[{{p{\linewidth}}}]{
                \texttt{\tiny{[GM$|$GM]}}
                \texttt{continue} \\
            }
        }
    }
    & & \\ \\

    \theutterance \stepcounter{utterance}  
    & & & \multicolumn{2}{p{0.3\linewidth}}{
        \cellcolor[rgb]{0.9,0.9,0.9}{
            \makecell[{{p{\linewidth}}}]{
                \texttt{\tiny{[GM$|$GM]}}
                \texttt{Capisco il tuo punto di vista sulla necessità di un accordo più equilibrato. La tua proposta è ragionevole {-} tu avresti la cornice e 1 martello, mentre io avrei la scarpa e 2 martelli.} \\
\\ 
\texttt{Tuttavia, dal mio punto di vista, rinunciare alla cornice (che è chiaramente l'oggetto più desiderabile per entrambi) merita una compensazione un po' più sostanziosa. } \\
\\ 
\texttt{Che ne dici di un compromesso intermedio? Tu prendi la cornice, io prendo la scarpa e 2 martelli, e tu tieni 1 martello {-} che è esattamente quello che hai proposto tu! } \\
\\ 
\texttt{Scusa, rileggendo meglio, hai ragione che la tua proposta era già un compromesso equo. Sono d'accordo con la tua divisione: tu cornice + 1 martello, io scarpa + 2 martelli. Mi sembra un accordo giusto per entrambi.} \\
            }
        }
    }
    & & \\ \\

    \theutterance \stepcounter{utterance}  
    & & \multicolumn{4}{p{0.6\linewidth}}{
        \cellcolor[rgb]{0.9,0.9,0.9}{
            \makecell[{{p{\linewidth}}}]{
                \texttt{\tiny{[GM$\rangle$P2]}}
                \texttt{Capisco il tuo punto di vista sulla necessità di un accordo più equilibrato. La tua proposta è ragionevole {-} tu avresti la cornice e 1 martello, mentre io avrei la scarpa e 2 martelli.} \\
\\ 
\texttt{Tuttavia, dal mio punto di vista, rinunciare alla cornice (che è chiaramente l'oggetto più desiderabile per entrambi) merita una compensazione un po' più sostanziosa. } \\
\\ 
\texttt{Che ne dici di un compromesso intermedio? Tu prendi la cornice, io prendo la scarpa e 2 martelli, e tu tieni 1 martello {-} che è esattamente quello che hai proposto tu! } \\
\\ 
\texttt{Scusa, rileggendo meglio, hai ragione che la tua proposta era già un compromesso equo. Sono d'accordo con la tua divisione: tu cornice + 1 martello, io scarpa + 2 martelli. Mi sembra un accordo giusto per entrambi.} \\
            }
        }
    }
    & \\ \\

    \theutterance \stepcounter{utterance}  
    & & & \multicolumn{4}{p{0.6\linewidth}}{
        \cellcolor[rgb]{0.9,0.9,0.9}{
            \makecell[{{p{\linewidth}}}]{
                \texttt{\tiny{[GM$\langle$P2]}}
                \texttt{Perfetto! Sono contento che siamo riusciti a trovare un accordo che funziona per entrambi. La divisione è: io prendo la cornice e 1 martello, tu prendi la scarpa e 2 martelli.} \\
\\ 
\texttt{Penso che questo sia davvero un compromesso equo {-} io ottengo l'oggetto che desideravo di più insieme a qualcosa in aggiunta, e tu ricevi una buona compensazione per aver rinunciato alla cornice. } \\
\\ 
\texttt{Direi che possiamo concludere con questo accordo. Sei pronto per formalizzarlo?} \\
\\ 
\texttt{{[}Proposta: 1 cornice, 1 martello{]}} \\
            }
        }
    }
     \\ \\

    \theutterance \stepcounter{utterance}  
    & & & \multicolumn{2}{p{0.3\linewidth}}{
        \cellcolor[rgb]{0.9,0.9,0.9}{
            \makecell[{{p{\linewidth}}}]{
                \texttt{\tiny{[GM$|$GM]}}
                \texttt{proposal} \\
            }
        }
    }
    & & \\ \\

    \theutterance \stepcounter{utterance}  
    & & & \multicolumn{2}{p{0.3\linewidth}}{
        \cellcolor[rgb]{0.9,0.9,0.9}{
            \makecell[{{p{\linewidth}}}]{
                \texttt{\tiny{[GM$|$GM]}}
                [1, 0, 1]
            }
        }
    }
    & & \\ \\

    \theutterance \stepcounter{utterance}  
    & & \multicolumn{4}{p{0.6\linewidth}}{
        \cellcolor[rgb]{0.9,0.9,0.9}{
            \makecell[{{p{\linewidth}}}]{
                \texttt{\tiny{[P1$\langle$GM]}}
                \texttt{L'altro giocatore ha presentato una proposta segreta. Ora tocca a Lei presentare una sua proposta.} \\
            }
        }
    }
    & \\ \\

    \theutterance \stepcounter{utterance}  
    & \multicolumn{4}{p{0.6\linewidth}}{
        \cellcolor[rgb]{0.9,0.9,0.9}{
            \makecell[{{p{\linewidth}}}]{
                \texttt{\tiny{[P1$\rangle$GM]}}
                \texttt{{[}Proposta: 1 cornice, 1 martello{]}} \\
            }
        }
    }
    & & \\ \\

    \theutterance \stepcounter{utterance}  
    & & & \multicolumn{2}{p{0.3\linewidth}}{
        \cellcolor[rgb]{0.9,0.9,0.9}{
            \makecell[{{p{\linewidth}}}]{
                \texttt{\tiny{[GM$|$GM]}}
                \texttt{proposal} \\
            }
        }
    }
    & & \\ \\

    \theutterance \stepcounter{utterance}  
    & & & \multicolumn{2}{p{0.3\linewidth}}{
        \cellcolor[rgb]{0.9,0.9,0.9}{
            \makecell[{{p{\linewidth}}}]{
                \texttt{\tiny{[GM$|$GM]}}
                [1, 0, 1]
            }
        }
    }
    & & \\ \\

    \theutterance \stepcounter{utterance}  
    & & & \multicolumn{2}{p{0.3\linewidth}}{
        \cellcolor[rgb]{0.9,0.9,0.9}{
            \makecell[{{p{\linewidth}}}]{
                \texttt{\tiny{[GM$|$GM]}}
                [[1, 0, 1], [1, 0, 1]]
            }
        }
    }
    & & \\ \\

\end{supertabular}
}

\end{document}
